\item Una lente delgada convergente con distancia focal de $14.0\text{ cm}$ forma una imagen del cuadrado $abcd$ de lado $10.0\text{ cm}$ ($h_c = h_b = 10.0\text{ cm}$) que se encuentra colocado sobre el eje óptico entre las distancias de $p_d = 20.0\text{ cm}$ y $p_a = 30.0\text{ cm}$. Sean $a', b', c'$ y $d'$ las esquinas correspondientes de la imagen.
\begin{enumerate}
    \item Determine las distancias de imagen $q_a$ (para $a'$ y $b'$), $q_d$ (para $c'$ y $d'$), y las alturas $h'_b$ y $h'_c$ de las esquinas superiores de la imagen.
    \item Elabore un bosquejo a escala de la forma geométrica distorsionada de la imagen.
    \item El área del objeto cuadrado es $100\text{ cm}^2$. Demuestre que para cualquier punto de la imagen a distancia $q$, la altura correspondiente $|h'|$ está dada por:
    $$ |h'| = \frac{10.0 q}{14.0 - q} $$
    donde las variables están en centímetros.
    \item Explique por qué el área geométrica de la imagen está dada por la integral:
    $$ A_{\text{imagen}} = \int_{q_a}^{q_d} |h'| dq $$
    \item Realice la integración correspondiente para calcular el área total de la imagen.
\end{enumerate}